\section{Related Work}
\label{sec-related}

\begin{table}[t]
\caption{Stores for ontologies, as documented (October 2026). We ran \sys, Oxigraph, Jena TDB2, and SQLite.}
\label{tab-related}
\small
\resizebox{\columnwidth}{!}{%
\begin{tabular}{@{}llll@{}}
\toprule
System & Triples are kept in & Hierarchy by & In a browser \\
\midrule
Oxigraph~\cite{oxigraph} & RocksDB (LSM) & query only & in-memory \\
Jena TDB2~\cite{jenatdb} & B+ trees & reasoner layer & no \\
RDF4J native~\cite{rdf4j} & B-trees & materialized & no \\
RDF-3X~\cite{neumann2010} & 6 compressed B+ trees & query only & no \\
RDFox~\cite{rdfox} & in memory & materialized & no \\
\sys & 3 B+ trees & intervals & full engine \\
\bottomrule
\end{tabular}}
\end{table}


\stitle{RDF stores}.
Hexastore indexes all six permutations~\cite{weiss2008}; RDF-3X keeps six
compressed B+ trees~\cite{neumann2010}; QLever builds read-mostly permutation
indexes~\cite{bast2017}; HDT is a compressed read-only format~\cite{fernandez2013};
Virtuoso stores quads in a relational engine~\cite{erling2010}.
Apache Jena TDB2 uses B+ trees~\cite{jenatdb}, RDF4J's native store
B-trees~\cite{rdf4j}, and Oxigraph~\cite{oxigraph} a log-structured merge tree in
RocksDB.
\sys follows the B+ tree line~\cite{bayer1972,comer1979} with three permutations;
its latching and recovery~\cite{lehman1981,graefe2011,mohan1992} are textbook, and
the contribution is the choice of parts, their measured behavior on ontologies,
and what is built on top.

\stitle{Inference}.
RDFS entailment is usually \emph{rewritten} into the query or \emph{materialized}
with maintenance by counting or delete-and-rederive~\cite{rdfox,gupta1993}.
Agrawal, Borgida, and Jagadish~\cite{agrawal1989} store a compressed transitive
closure by numbering a spanning forest in postorder with a set of intervals per
node, and Grust~\cite{grust2002} numbers XML trees so that axes become regions of
a plane; we use the former on class DAGs inside a B+ tree store and make it the
target of a plan rewrite.

\stitle{Vector search}.
Approximate nearest neighbor search has inverted files and product
quantization~\cite{jegou2011}, proximity graphs such as HNSW~\cite{malkov2020},
and disk-oriented graphs such as DiskANN~\cite{subramanya2019}; FAISS packages
the first two~\cite{johnson2021}, and Pan, Wang, and Li survey the hybrid queries
that are hard for them~\cite{pan2024}; ACORN makes HNSW traverse only nodes that
pass a predicate~\cite{patel2024}.
Neo4j, Neptune Analytics, Fluree, and GraphDB add vector or similarity indexes to
graph and RDF stores~\cite{neo4jvec,neptunevec,fluree,graphdbsim}; in the
documentation we read, none restricts the search with range-scannable hierarchy
intervals in the engine that holds the triples.

\stitle{Adaptive radix trees and browsers}.
An ART~\cite{leis2013} orders keys by byte, so its range scan is lexical and
Lemma~\ref{lem-bw} applies; sibling links are awkward there, because lazy
expansion moves a leaf, and \sys's three fixed-width integers leave an ART little
to compress.
WebAssembly~\cite{haas2017} lets a C++ engine run in a page; Oxigraph ships an
in-memory JavaScript build, and \sys compiles the whole paged engine.
